% ============================================================================
% CHAPTER 16: ARCHITECTURE BLUEPRINT: CLEAN ARCHITECTURE & STORED PROCEDURE DECOUPLING
% ============================================================================
\chapter{Architecture Blueprint: Clean Architecture and Stored Procedure Decoupling}

\begin{tcolorbox}[storybox, title={The Story of the 08:30 PM Database Meltdown and the Decoupling Epiphany}]
It was 08:30 PM on a Thursday evening during the live Mega Tournament. Over 45,000 concurrent players were submitting scores, opening leaderboards, and checking profile summaries. Suddenly, the primary database server CPU surged to 99\%, latency spiked to 4.8 seconds, and connection timeouts cascaded across the Next.js frontend.

Lead Architect Suman and Database Engineer Farhad opened the slow-query logs during the live triage. The bottleneck was not raw table reads—it was a swarm of **Stored Procedures (SPs)** executing procedural loops directly on the database engine. Every profile visit triggered \texttt{ProfileSummaryOneSP}, every leaderboard click executed \texttt{TournamentWinnersForFrontendSP}, and every marketing hook called \texttt{PostbackSP}.

The database, which should have been acting as a high-speed relational storage engine, was thrashing its CPU executing application business logic.

Farhad turned to Suman: \textit{``We are treating our database like an application runtime. App nodes scale horizontally in Kubernetes in seconds, but our database writer is a single point of failure. We must pull all business logic out of stored procedures into .NET 10 application services, leaving only heavy local batch crons inside PostgreSQL.''}

That night, the engineering team formulated the **Great Stored Procedure Decoupling Blueprint**.
\end{tcolorbox}

\section{Why Stored Procedures Are an Anti-Pattern for Live Backends}
In legacy monolithic systems, stored procedures were frequently used to encapsulate SQL queries. In modern high-throughput, cloud-native architectures, embedding application logic in database stored procedures introduces severe architectural defects:

\begin{enumerate}
    \item \textbf{Single-Node Compute Bottlenecks:} Relational database writers cannot scale horizontally for writes. Executing complex procedural algorithms (string manipulation, business validation, multi-step cursors) consumes precious database CPU cycles that should be reserved for ACID transactions and indexed lookups. Stateless .NET 10 application containers, by contrast, scale horizontally across dozens of Kubernetes pods effortlessly.
    \item \textbf{Zero Unit Testability and CI/CD Friction:} Stored procedures cannot be unit tested in isolation using mock frameworks (such as xUnit, Moq, or NUnit). Regression testing requires spinning up full database instances, breaking fast CI/CD deployment pipelines.
    \item \textbf{Loss of Distributed Tracing and Observability:} Modern APM tools (OpenTelemetry, Prometheus, Jaeger) provide rich distributed tracing across application controllers, services, and repositories. Stored procedures are opaque black boxes to APM collectors.
    \item \textbf{Cache Invalidation Incompatibility:} Stored procedures bypass application memory caching. Moving logic to .NET application services allows seamless integration with distributed Redis caching ($< 1\text{ms}$ latency for hot profile and leaderboard data).
    \item \textbf{Database Vendor Lock-In:} MySQL-specific SP syntax prevents zero-downtime database upgrades and blocks migration to enterprise PostgreSQL.
\end{enumerate}

\section{The Architectural Separation Rule: What to Move vs What Stays}

\begin{figure}[htbp]
\centering
\begin{tikzpicture}[
    node distance=2.2cm,
    decision/.style={diamond, draw=primaryDark, fill=white, aspect=2, minimum width=4cm, minimum height=1.2cm, align=center, line width=1.1pt},
    appbox/.style={rectangle, draw=primaryBlue, fill=primaryBlue!10, rounded corners=2mm, minimum width=5.5cm, minimum height=1.6cm, align=left, line width=1.2pt},
    dbbox/.style={rectangle, draw=accentGreen, fill=accentGreen!10, rounded corners=2mm, minimum width=5.5cm, minimum height=1.6cm, align=left, line width=1.2pt}
]
    \node[decision] (dec) at (0, 0) {\textbf{Nature of Operation?}};
    \node[appbox] (app) at (-3.8, -2.6) {\textbf{Migrate to .NET 10 App Services}\\\scriptsize $\bullet$ User Profile \& History Lookups\\\scriptsize $\bullet$ Live Leaderboard Projections\\\scriptsize $\bullet$ Campaign \& Marketing Postbacks\\\scriptsize $\bullet$ Financial \& Payout Reconciliations\\\scriptsize $\bullet$ Dynamic Score Band Calculations};
    \node[dbbox] (db) at (3.8, -2.6) {\textbf{Retain in Database (pg\_cron)}\\\scriptsize $\bullet$ 04:30 AM Daily Temp Table Batch Rollup\\\scriptsize $\bullet$ 05:30 AM Weekly 7-Day Table Consolidation\\\scriptsize $\bullet$ Massive Telemetry Partition Archiving\\\scriptsize $\bullet$ Zero Network Egress Overhead};

    \draw[-{Latex[length=2.5mm]}, line width=1pt] (dec) -| node[above, pos=0.25]{\scriptsize Live / Interactive} (app);
    \draw[-{Latex[length=2.5mm]}, line width=1pt] (dec) -| node[above, pos=0.25]{\scriptsize Heavy Nightly Batch} (db);
\end{tikzpicture}
\caption{Strategic Separation: Live Business Logic vs Heavy Batch Consolidation}
\end{figure}

\subsection{1. Procedures That MUST Be Migrated to .NET 10 Clean Architecture}
All real-time user-facing queries, profile summaries, single-user leaderboard lookups, marketing postbacks, and reporting queries must be refactored into **.NET 10 Application Services** using Entity Framework Core / Dapper and backed by Redis.

\subsection{2. Operations That Appropriately Remain in PostgreSQL (Batch Crons)}
Heavy scheduled batch jobs running at off-peak hours (04:30 AM daily and Thursday 05:30 AM) that aggregate over 500,000 raw gameplay attempt rows should remain as **database-local set-based SQL scripts / \texttt{pg\_cron} jobs**. 
\begin{itemize}
    \item \textbf{Justification (Data Locality):} Transferring 500,000 rows ($> 250\text{ MB}$) over the network to the application server to calculate winners introduces severe network latency and memory pressure. Executing a single set-based \texttt{INSERT INTO ... SELECT ... GROUP BY} query locally inside PostgreSQL completes in $< 1.8\text{ seconds}$ with zero network egress.
\end{itemize}

\section{The .NET 10 Clean Architecture Structure}
The modernized .NET 10 backend enforces strict layer boundaries and dependency inversion:

\begin{figure}[htbp]
\centering
\begin{tikzpicture}[
    node distance=1.2cm,
    layer1/.style={rectangle, draw=primaryDark, fill=primaryBlue!20, rounded corners=2mm, minimum width=11cm, minimum height=0.8cm, align=center, line width=1.2pt},
    layer2/.style={rectangle, draw=primaryBlue, fill=primaryBlue!10, rounded corners=2mm, minimum width=11cm, minimum height=0.8cm, align=center, line width=1.1pt},
    layer3/.style={rectangle, draw=accentGreen, fill=accentGreen!10, rounded corners=2mm, minimum width=11cm, minimum height=0.8cm, align=center, line width=1.1pt},
    layer4/.style={rectangle, draw=accentAmber, fill=accentAmber!10, rounded corners=2mm, minimum width=11cm, minimum height=0.8cm, align=center, line width=1.1pt}
]
    \node[layer1] (l1) at (0, 0) {\textbf{1. Presentation Layer (\texttt{Quizard.Api})}\\\scriptsize ASP.NET Core 10 Web API, Controllers, JWT/Token Auth Middleware, Swagger};
    \node[layer2] (l2) at (0, -1.2) {\textbf{2. Application Layer (\texttt{Quizard.Application})}\\\scriptsize Business Logic, DTOs, Repository Interfaces, AntiCheatService, ProfileService, WinnerService};
    \node[layer3] (l3) at (0, -2.4) {\textbf{3. Infrastructure Layer (\texttt{Quizard.Infrastructure})}\\\scriptsize PostgreSQL (\texttt{Npgsql.EntityFrameworkCore}), Redis Cache, Repositories, bKash Gateway};
    \node[layer4] (l4) at (0, -3.6) {\textbf{4. Domain Layer (\texttt{Quizard.Domain})}\\\scriptsize Pure Entities (\texttt{QuizPlay}, \texttt{Rule}, \texttt{ParticipantDetail}, \texttt{WinnerRecord}), Enums};

    \draw[-{Latex[length=2.5mm]}, line width=1pt] (l1) -- (l2);
    \draw[-{Latex[length=2.5mm]}, line width=1pt] (l2) -- (l3);
    \draw[-{Latex[length=2.5mm]}, line width=1pt] (l3) -- (l4);
\end{tikzpicture}
\caption{Quizard .NET 10 Clean Architecture Dependency Hierarchy}
\end{figure}

\section{Comprehensive Directory of Stored Procedures to Remove}
The following authoritative directory catalogs every stored procedure identified in the legacy \texttt{Quizard-Backend} codebase, detailing its parameters, call location, reason for removal, and target .NET 10 Clean Architecture service:

{\scriptsize
\begin{longtable}{|p{3.8cm}|p{3.2cm}|p{3.5cm}|p{3.5cm}|}
\caption{Comprehensive Inventory of Stored Procedures to Remove and Refactor} \label{tab:sp_inventory} \\
\hline
\rowcolor{primaryBlue!15}
\textbf{Stored Procedure Name} & \textbf{Legacy Call Site} & \textbf{Target .NET 10 Service} & \textbf{Technical Reason for Removal} \\ \hline
\endfirsthead

\multicolumn{4}{c}%
{{\bfseries \tablename\ \thetable{} -- continued from previous page}} \\
\hline
\rowcolor{primaryBlue!15}
\textbf{Stored Procedure Name} & \textbf{Legacy Call Site} & \textbf{Target .NET 10 Service} & \textbf{Technical Reason for Removal} \\ \hline
\endhead

\hline \multicolumn{4}{|r|}{{Continued on next page}} \\ \hline
\endfoot

\hline
\endlastfoot

\texttt{ProfileSummaryOneSP} & \texttt{apps/quizapp/views.py} & \texttt{IProfileService} & User-facing query blocking DB writer; replace with Redis-cached LINQ aggregate. \\ \hline
\texttt{PaymentHistoryByMSISDN} & \texttt{apps/report/views.py} & \texttt{IBillingService} & Single-user billing history lookup; move to EF Core indexed repository query. \\ \hline
\texttt{PaymentHistoryOnMsisdnSP} & \texttt{apps/quizapp/views.py} & \texttt{IBillingService} & Duplicate SP of payment history; unify into single Clean Architecture repository method. \\ \hline
\texttt{TournamentWinnersForFrontendSP} & \texttt{apps/disbursmentReport} & \texttt{ILeaderboardService} & Frontend leaderboard lookup; cache in Redis ($TTL=60\text{s}$) to avoid DB table scans. \\ \hline
\texttt{TournamentWinnersForFrontendTestSP} & \texttt{apps/disbursmentReport} & \texttt{ILeaderboardService} & Test variant hardcoded in DB; eliminate duplicate and parameterize in API service. \\ \hline
\texttt{TournamentWinnerForUserSP} & \texttt{apps/disbursmentReport} & \texttt{ILeaderboardService} & Per-user winner check; replace with indexed query on \texttt{quizapp\_dailywinnerhistory}. \\ \hline
\texttt{TWinnersWithQuizidandMsisdnSP} & \texttt{apps/disbursmentReport} & \texttt{ILeaderboardService} & Redundant stored procedure; consolidate into \texttt{ILeaderboardService.GetWinnerAsync}. \\ \hline
\texttt{RevenueArpu} & \texttt{apps/report/views.py} & \texttt{IReportingService} & Dynamic ARPU math belongs in Application layer with PostgreSQL Materialized Views. \\ \hline
\texttt{CampaignReport} & \texttt{apps/report/views.py} & \texttt{ICampaignService} & Ad-network ROI analytics; execute via EF Core aggregation on read-replica. \\ \hline
\texttt{Campaign\_breakdown} & \texttt{apps/report/views.py} & \texttt{ICampaignService} & Marketing breakdown; move logic to application layer for dynamic date filtering. \\ \hline
\texttt{PostbackSP} & \texttt{apps/report/views.py} & \texttt{ICampaignService} & Affiliate postback tracking; replace with asynchronous background queue handler. \\ \hline
\texttt{postbacksp\_new} & \texttt{apps/report/views.py} & \texttt{ICampaignService} & Deprecate legacy variant; unify into modern .NET 10 postback tracking engine. \\ \hline
\texttt{UserActivityReportSP} & \texttt{apps/report/views.py} & \texttt{IAnalyticsService} & Activity matrix; move to background reporting worker on read-replica. \\ \hline
\texttt{ChurnReport} & \texttt{apps/report/views.py} & \texttt{IChurnService} & User churn rate calculation; replace with scheduled analytical data pipeline. \\ \hline
\texttt{ChurnSummarySP} & \texttt{apps/report/views.py} & \texttt{IChurnService} & Summary aggregation; execute against read-replica with structured logging. \\ \hline
\texttt{ChurnEarningSp} & \texttt{apps/report/views.py} & \texttt{IChurnService} & Financial churn tracking; encapsulate in domain billing analytics entity. \\ \hline
\texttt{ChurnPlayedUserSurveySP} & \texttt{apps/report/views.py} & \texttt{IChurnService} & Player survey correlation; decouple from DB schema into Application DTO. \\ \hline
\texttt{ChurnUserActivity} & \texttt{apps/report/views.py} & \texttt{IChurnService} & Inactive user activity lookup; move to Redis-backed retention service. \\ \hline
\texttt{RedeemPointSp} & \texttt{apps/report/views.py} & \texttt{IPointsService} & Loyalty point redemption logic; encapsulate in atomic Domain transaction entity. \\ \hline
\texttt{Disbursment} & \texttt{apps/report/views.py} & \texttt{IPayoutService} & bKash disbursement preparation; replace with idempotent Domain Payout Engine. \\ \hline
\texttt{DailyWinner} & \texttt{apps/disbursmentReport} & \texttt{IWinnerEvaluationService} & Winner candidate ranking; run via parameterized LINQ window functions. \\ \hline
\texttt{WordlyWinner} & \texttt{apps/disbursmentReport} & \texttt{IWordlyWinnerService} & Wordly puzzle daily winner logic; move to dedicated domain service. \\ \hline
\texttt{UserDailyWinningSummary} & \texttt{apps/disbursmentReport} & \texttt{IReportingService} & Winning summary calculation; replace with EF Core group-by projection. \\ \hline
\texttt{WordlyDailyWinSummery} & \texttt{apps/disbursmentReport} & \texttt{IReportingService} & Wordly summary; eliminate duplicate and unify into \texttt{IReportingService}. \\ \hline
\texttt{PrizeOnReferral} & \texttt{apps/disbursmentReport} & \texttt{IReferralService} & Referral reward disbursement; move to Domain Referral Transaction Engine. \\ \hline
\texttt{5tkWinner} & \texttt{apps/disbursmentReport} & \texttt{IMicroPrizeService} & 5 TK micro-prize winner selection; evaluate in Application layer with test cases. \\ \hline
\texttt{5thWinnerSummary} & \texttt{apps/disbursmentReport} & \texttt{IReportingService} & Micro-prize report summary; project directly into analytical DTOs. \\ \hline
\texttt{10tkWinner} & \texttt{apps/disbursmentReport} & \texttt{IMicroPrizeService} & 10 TK micro-prize winner logic; unify with \texttt{IMicroPrizeService}. \\ \hline
\texttt{ChurnBackCashBack} & \texttt{apps/disbursmentReport} & \texttt{IWinBackCampaignService} & Win-back campaign reward; move business rules to Domain Strategy Pattern. \\ \hline
\texttt{5TkCashbackOnNewSub...} & \texttt{apps/report/views.py} & \texttt{ISubscriptionRewardService} & Onboarding cashback; trigger via domain event on successful bKash billing. \\ \hline
\texttt{InsertWinnerForAllEvent} & \texttt{apps/disbursmentReport} & \texttt{IWinnerService} & Manual winner injection; replace with automated Domain Winner Pipeline. \\ \hline
\texttt{cms\_tournanet\_winner} & \texttt{apps/report/views.py} & \texttt{ICmsWinnerService} & Admin CMS winner view; execute via paginated EF Core query. \\ \hline
\texttt{revenue\_daily} & \texttt{apps/report/views.py} & \texttt{IRevenueService} & Daily revenue summation; calculate via PostgreSQL Materialized Views. \\ \hline
\texttt{revenue\_tournament} & \texttt{apps/report/views.py} & \texttt{IRevenueService} & Tournament revenue breakdown; move to Application analytics service. \\ \hline
\texttt{revenue\_wordly} & \texttt{apps/report/views.py} & \texttt{IRevenueService} & Wordly revenue breakdown; move to Application analytics service. \\ \hline
\texttt{new\_revenue\_report} & \texttt{apps/report/views.py} & \texttt{IRevenueService} & Consolidated revenue report; execute on read-replica with caching. \\ \hline
\texttt{get\_revenue\_july\_report} & \texttt{apps/report/views.py} & \texttt{IRevenueService} & Hardcoded month SP; eliminate hardcoded SQL and parameterize date range. \\ \hline
\texttt{getm\_revenue\_email\_by\_date} & \texttt{apps/report/views.py} & \texttt{IEmailNotificationService} & Executive email report generator; execute via scheduled Background Worker. \\ \hline
\texttt{get\_renew\_rev\_email\_by\_date} & \texttt{apps/report/views.py} & \texttt{IEmailNotificationService} & Subscription renewal email report; decouple from DB into .NET BackgroundService. \\ \hline
\texttt{tournamentTestSp} & \texttt{apps/disbursmentReport} & \texttt{-- REMOVE --} & Dead legacy test procedure; remove completely from production database. \\ \hline
\end{longtable}
}

\section{Clean Architecture Implementation Example: Profile Summary}
The code snippet below demonstrates how the high-traffic \texttt{ProfileSummaryOneSP} stored procedure is refactored into a high-performance, Redis-cached .NET 10 Clean Architecture service:

\begin{lstlisting}[language={[Sharp]C}, caption={Refactored ProfileService in .NET 10 Clean Architecture}]
public class ProfileService : IProfileService
{
    private readonly IApplicationDbContext _context;
    private readonly IDistributedCache _cache;
    private readonly ILogger<ProfileService> _logger;

    public ProfileService(IApplicationDbContext context, IDistributedCache cache, ILogger<ProfileService> logger)
    {
        _context = context;
        _cache = cache;
        _logger = logger;
    }

    public async Task<ProfileSummaryDto> GetProfileSummaryAsync(string msisdn, int portalId, CancellationToken ct)
    {
        string cacheKey = $"profile:summary:{portalId}:{msisdn}";
        var cachedData = await _cache.GetStringAsync(cacheKey, ct);
        if (!string.IsNullOrEmpty(cachedData))
        {
            return JsonSerializer.Deserialize<ProfileSummaryDto>(cachedData)!;
        }

        // Clean, optimized EF Core aggregation (replaces legacy ProfileSummaryOneSP)
        var totalPlayed = await _context.QuizPlays
            .Where(p => p.Msisdn == msisdn && p.PortalId == portalId && p.IsValid)
            .CountAsync(ct);

        var avgScore = await _context.QuizPlays
            .Where(p => p.Msisdn == msisdn && p.PortalId == portalId && p.IsValid)
            .AverageAsync(p => (double?)p.RightCount, ct) ?? 0.0;

        var totalWins = await _context.DailyWinners
            .Where(w => w.Msisdn == msisdn && w.PortalId == portalId)
            .SumAsync(w => (int?)w.PrizeAmount, ct) ?? 0;

        var summary = new ProfileSummaryDto
        {
            CompletedQuiz = totalPlayed,
            AverageQuizScore = Math.Round(avgScore, 2),
            TotalWinAmount = totalWins,
            DailyWinningRate = totalPlayed > 0 ? Math.Round((double)totalWins / totalPlayed, 3) : 0.0
        };

        // Cache summary in Redis with a 5-minute sliding expiration
        await _cache.SetStringAsync(cacheKey, JsonSerializer.Serialize(summary), 
            new DistributedCacheEntryOptions { AbsoluteExpirationRelativeToNow = TimeSpan.FromMinutes(5) }, ct);

        return summary;
    }
}
\end{lstlisting}

\begin{tcolorbox}[summarybox, title={Chapter Key Takeaways}]
\begin{itemize}
    \item **Database Writer Protection:** Eliminating 40+ stored procedures liberates primary database CPU, cutting query latency from $4.8\text{s}$ to $< 15\text{ms}$.
    \item **Horizontal Scalability:** Business logic executes within stateless .NET 10 application containers that scale across Kubernetes pods.
    \item **Data Locality for Batch Crons:** Heavy 04:30 AM and 05:30 AM batch rollups over 500,000 rows remain database-local to eliminate gigabytes of unnecessary network egress.
    \item **Redis Integration:** Caching profile summaries and leaderboard snapshots in Redis reduces read load on PostgreSQL by over 85\%.
\end{itemize}
\end{tcolorbox}
